\section{Background and Related Work}
\label{sec:background}

\subsection{AI-Like Terminal Interaction Model}
\label{sec:ai-like-model}

AI-like terminal interaction is organized around a closed feedback loop. A controller issues an action, observes textual or visual output, and then selects the next action based on that observation \cite{yao2023react}. SWE-agent exposes this pattern through an agent-computer interface in which software-engineering agents use shell commands, file operations, and terminal observations to make progress on tasks \cite{Yang2024}. Terminal-Bench similarly treats command-line environments as the task substrate for evaluating whether agents can complete realistic terminal tasks \cite{merrill2026}. In such a loop, the remote shell is not only a login mechanism; it is the observation channel through which the client learns whether a command completed, which data were produced, and when an interactive edit became visible.

This study abstracts that terminal-facing channel rather than executing a full agent. The benchmark client deterministically issues commands or probes and records client-observable events. This design removes model reasoning quality, planning errors, and task-specific policies from the experiment, allowing the measurement to focus on transport behavior. The abstraction is intentionally narrower than an end-to-end agent benchmark, but it exposes the protocol properties that such benchmarks normally assume: low setup latency, timely feedback, and complete parseable output.

\subsection{Remote Shell Protocols}
\label{subsec:protocols}

Remote shell protocols differ along three performance-relevant dimensions for AI-like terminal clients: how quickly a session becomes usable, how packet loss affects visible feedback, and whether the client receives a complete byte stream or only the latest terminal state.

\paragraph{SSHv2}
The Secure Shell Protocol version 2 \cite{Ylonen2006RFC4251} is the de facto mechanism for secure remote terminal access. It runs over TCP and provides ordered, reliable delivery for encrypted SSH channels \cite{Bellare2002}. This byte-stream semantics is useful when the client must receive complete command output, but it also inherits TCP HOL blocking: a lost segment can delay later data until retransmission completes \cite{Scharf2007}. For a closed-loop terminal client, this can appear as delayed command completion or delayed visible feedback. SSHv2 setup also includes a TCP handshake followed by SSH key exchange and authentication, so the setup cost grows with RTT \cite{Ylonen2006RFC4253}.

\paragraph{SSH3}
SSH3 redesigns SSH by carrying secure-shell semantics over HTTP/3 \cite{michel2023}. HTTP/3 relies on QUIC, which provides encrypted transport, connection establishment integrated with TLS, and multiplexed streams over UDP \cite{Iyengar2021}. Its main architectural promise is lower setup latency: QUIC integrates transport and cryptographic negotiation, enabling 1-RTT setup and 0-RTT resumption for eligible connections. QUIC also removes HOL blocking across independent streams. For terminal workloads, however, the benefit depends on how the shell traffic is mapped onto streams and on whether the workload contains independent streams or a single ordered output flow. A large command printed through one terminal stream can still be limited by ordered delivery and retransmission within that stream, even though unrelated QUIC streams would not be blocked.

\paragraph{Mosh}
Mosh makes the opposite trade-off \cite{Winstein2012}. It uses SSH only to authenticate and bootstrap the session, then switches to a UDP channel carrying terminal state updates through the State Synchronization Protocol (SSP). SSP treats the terminal as replicated state and sends compact updates for the latest screen. This design avoids waiting for every intermediate byte and makes lost datagrams superseded by newer state. Mosh further improves perceived responsiveness through predictive local echo. The cost is completeness: output that scrolls past between screen updates may never be delivered to the client. This is acceptable for many human editing sessions, but it is risky for automated clients that parse command output as data.

\paragraph{Prior Measurement Studies}
Prior work provides useful but fragmented evidence. Mosh was originally evaluated for human-perceived responsiveness under mobile connectivity, showing that state synchronization and local echo can substantially reduce visible keystroke delay compared with SSHv2 \cite{Winstein2012}. Google's deployment study presents QUIC as a transport that combines secure connection setup, multiplexing, and congestion-control changes to improve selected web workloads at Internet scale \cite{Langley2017}. A later measurement study cautions that QUIC results depend on implementation version, workload, and network path, so comparisons against TCP need controlled methodology rather than assuming a uniform QUIC advantage \cite{Kakhki2017}. SWE-agent evaluates whether an agent can solve software-engineering tasks through a command-line interface and related tool abstractions \cite{Yang2024}. Terminal-Bench evaluates whether agents can complete realistic tasks inside terminal environments \cite{merrill2026}. These benchmarks stress task success and agent behavior, but they do not separately measure the remote-shell transport events that this paper studies: session setup, command-completion latency, keystroke visibility, and byte-complete output delivery.

These prior results point to distinct terminal requirements rather than one universal performance objective. A short command loop is sensitive to request-response delay after a session is established. A large command output requires the client to receive complete and content-correct bytes, not merely the latest visible screen. Interactive editing depends on when a character becomes visible to the client, which may benefit from prediction or state synchronization even when byte-level completeness is not preserved. What remains missing is a controlled comparison of SSHv2, Mosh, and SSH3 on these separate terminal-event workloads, with latency and output completeness measured under the same network scenarios. This paper fills that gap for AI-like interactive terminal clients.
